\section{导读：为什么这几份报告都指向 Agent}

本节先建立整期的学习目标。这期节目读了几份近期最值得细读的 Agent 技术材料：Kimi K2 技术报告、ChatGPT Agent 发布文、Qwen3-Coder 技术博文，以及 Manus 的上下文工程经验。它们看似来自不同公司和产品，但共同指向一个主题：大模型正在从“回答问题”走向“在环境中完成任务”。

Agent 的关键不只是调用工具，而是完整闭环：感知环境、决策、行动、读取反馈、修正轨迹。这个闭环会牵出一整套工程问题：合成数据、轨迹生成、强化学习、可验证奖励、安全沙盒、上下文工程、KV cache、工具选择、文件系统记忆和多 Agent 组织。

\begin{figure}[H]
\centering
\includegraphics[width=0.96\textwidth]{report-comparison.png}
\caption{四份报告对照：Kimi、ChatGPT Agent、Qwen3-Coder、Manus 分别强调不同层。自制概念图，依据视频结构和官方材料整理。}
\end{figure}

\begin{knowledgebox}{读图：四份材料各看一层}
Kimi K2 看开放模型和 agentic training，ChatGPT Agent 看统一产品系统，Qwen3-Coder 看代码环境里的 RL 和工具使用，Manus 看生产 Agent 的上下文工程。合起来就是 Agent 从模型、训练、产品到工程部署的全链条。
\end{knowledgebox}

\begin{importantbox}{本期核心命题}
Agent 的进步不是单一模型能力提升，而是系统工程：模型、环境、任务、工具、奖励、安全和上下文共同优化，才能把智能从输出文本推进到完成真实任务。
\end{importantbox}

\begin{knowledgebox}{四份材料的学习顺序}
\begin{tabularx}{\textwidth}{p{0.20\textwidth}p{0.34\textwidth}X}
\toprule
材料 & 先看什么 & 教学价值 \\
\midrule
Kimi K2 & MoE、MuonClip、agentic data、joint RL & 看开放模型如何被训练成 Agentic。 \\
ChatGPT Agent & Operator + Deep Research + 工具计算机 & 看 Agent 如何变成用户产品。 \\
Qwen3-Coder & Code RL、long-horizon RL、Qwen Code & 看代码环境为什么适合 Agent RL。 \\
Manus & KV cache、工具遮蔽、文件系统、错误保留 & 看生产 Agent 如何靠上下文工程稳定运行。 \\
\bottomrule
\end{tabularx}
\end{knowledgebox}

\begin{knowledgebox}{关键数字速览}
\begin{tabularx}{\textwidth}{p{0.22\textwidth}p{0.34\textwidth}X}
\toprule
材料 & 关键数字/事实 & 教学含义 \\
\midrule
Kimi K2 & 约 1T 总参数、32B 激活参数、15.5T tokens & 开放 MoE 模型也进入超大规模系统训练。 \\
Kimi K2 & SWE-Bench Verified 65.8、Tau2-Bench 66.1 & Agentic benchmark 成为模型报告核心指标。 \\
Qwen3-Coder & 480B MoE、35B active、256K/1M context & 代码 Agent 需要超长上下文和仓库级理解。 \\
Qwen3-Coder & 20,000 并行环境用于 long-horizon RL & Agent RL 的瓶颈是环境规模化。 \\
Manus & 输入/输出 token 比例约 100:1 & KV cache 对成本和延迟极其关键。 \\
\bottomrule
\end{tabularx}
\end{knowledgebox}

\subsection{本章小结}

EP110 是 Agent 技术报告课。它把 EP115 的 Agent 理论、EP113 的 K2 模型公司视角、EP116 的企业级 Agentic Model 连接起来，展示“系统工程的力量”。

